\section{PVL-Delta Model}

\PVL{} first scales the objective net outcome in the same way:
\begin{equation}
    x_t=\frac{\mathrm{win}_t+\mathrm{loss}_t}{100}.
\end{equation}
It then transforms that outcome into subjective utility:
\begin{equation}
    u(x_t)=
    \begin{cases}
        x_t^{\alpha},            & x_t\ge 0, \\
        -\lambda(-x_t)^{\alpha}, & x_t<0.
    \end{cases}
\end{equation}
The outcome-sensitivity parameter \(\alpha\) controls the nonlinear effect of outcome magnitude, while the loss-aversion parameter \(\lambda\) controls the relative weighting of negative outcomes.

The subjective utility becomes the teaching signal in the delta update:
\begin{equation}
    \delta_t=u(x_t)-E_{a_t}(t),
\end{equation}
\begin{equation}
    E_{a_t}(t+1)=E_{a_t}(t)+\phi\delta_t,
\end{equation}
where \(\phi\) is the learning rate and only the chosen deck expectancy is updated.

The fitted response-consistency parameter \(c\) is transformed into choice sensitivity as
\begin{equation}
    \theta=3^c-1.
\end{equation}
Thus \(c\) is fitted, whereas \(\theta\) is derived. Choice probabilities are
\begin{equation}
    P(a_t=d)=\frac{\exp\left[\theta E_d(t)\right]}{\sum_{j=1}^{4}\exp\left[\theta E_j(t)\right]}.
\end{equation}
The four fitted parameters are therefore \(\phi\), \(\alpha\), \(\lambda\), and \(c\)~\cite{ahn2008}.
